\chapter{Implementing Horn Clause Logic on RED2}

The eventual goal of the RED project is to implement a high-performance
computing system that symmetrically combines the complete
\(\lambda\beta\eta\)-calculus and the full first order predicate calculus. At
the present, this goal is out of reach. Neither \(\eta\)-reduction nor full
first order logic can currently be implemented in a way that could be called
``high-performance''. This situation may change in the near future; recent work
by Stickel~[58] and Loveland~[37] indicates that efficient implementation of
full first order logic may not be far off. At the present time, however, we
must be satisfied to limit ourselves to the \lambdabetacalculus{} and the Horn
clause subset of first order logic.

I shall conclude this dissertation with a discussion of an experiment in
implementing Horn clause logic on RED2. The approach taken was to explore the
minimal architectural changes necessary to integrate some form of Horn clause
logic programming capability into THOR which conforms to THOR's incremental
reduction semantics. The result is THOR-L, a language where relations are
represented as boolean functions and a single computed answer substitution is
made available through the environment. THOR-L is not a complete success from a
logic programming standpoint because it is unable to make available (in a truly
useful way) multiple answers resulting from nondeterministic relations.

This chapter assumes the reader is familiar with basic logic programming
concepts and terminology; the necessary background information can be found
in~[36].

\section{Logic Programming}

A logic program is a set of definite Horn clauses. A logic program is executed
by performing resolution~[48] on the set of program clauses and a goal clause
containing existentially quantified variables. Resolution is a nondeterministic
rewriting process which determines if the goal is a logical consequence of the
program. In the process of determining logical consequence, resolution may
compute bindings for some (or all) of the existential variables in the goal. In
general, it is the values of these bindings which are of interest; if the goal
is a logical consequence of the program, these bindings are called a correct
answer substitution. Because relations may be nondeterministic, there may be
more than one correct answer substitution.

The resolution procedure has two basic operations: unification and search.
Unification attempts to make a predicate in the goal equal to a predicate which
is the consequent of a program clause by instantiating variables. When equality
can be achieved, the antecedent of the program clause replaces the matched
predicate in the goal and the bindings generated by unification are applied to
this new goal clause. Search is the process of selecting which program clause to
resolve upon; at any given time, there may be several applicable program
clauses.

\section{Integrating THOR and Logic Programming}

At first glance, it may seem that \lambdacalculus{} systems based on reduction
and relational systems based on resolution are fundamentally different. This is
not the case; it is possible to design functional languages which are executed
by techniques similar to resolution [46] [45] [18] [8], and relational languages
which are executed by reduction [49] [10] [17].

The RED project seeks to combine both functions and relations symmetrically in a
single coherent framework based upon reduction.\footnote{A symmetrical
combination allows relations to be used in functions and vice versa.} As a step
towards this goal, an attempt was made to integrate THOR with Horn clause logic
programming on the RED2 architecture. The resulting combination, THOR-L, was
shaped by the following goals:

\begin{enumerate}
\item \emph{Notation} The integrated language should provide a uniform syntax
for both functions and relations, and this language should compile to a
RED2-style instruction set that supports de-compilation.
\item \emph{Unification} A unification primitive must be added to RED2 to
support resolution. This primitive should conform as much as possible to the
existing memory regimen and incremental reduction semantics of RED2.
\item \emph{Search} A search strategy for resolution must be devised which
requires little, if any, modification to the existing RED2 architecture.
\item \emph{Language Interface} The results of functions and relations should
be made available to one another in a natural way.
\end{enumerate}

None of these goals is independent of the others; they all interact in a
synergistic way. The design and implementation of THOR-L is fairly successful
with regard to goals 1, 3, and 4. A portion of the second goal, incorporating
unification without altering the RED2 architecture, proved impossible to
achieve.

\subsection{A Uniform Notation}

Each function in THOR is expressed as single definition. In most logic
programming languages each relation is expressed as a set of clauses, but they
can also be expressed as a single definition by collecting all of their clauses
together in a manner suggested by Clark~[14]. Suppose that
\[
R(t_1,\ldots,t_n)\leftarrow L_1\wedge\cdots\wedge L_m
\tag{6.1}
\]
is a clause in relation \(R\). If \(z_1,\ldots,z_n\) are variables not appearing
in the clause, then (6.1) is equivalent to the clause
\[
R(z_1,\ldots,z_n)\leftarrow z_1=t_1\wedge\cdots\wedge z_n=t_n
\wedge L_1\wedge\cdots\wedge L_m.
\]
If \(y_1,\ldots,y_p\) are the variables of (6.1), the previous equation now
becomes
\[
R(z_1,\ldots,z_n)\leftarrow
(\exists y_1\ldots y_p)[z_1=t_1\wedge\cdots\wedge z_n=t_n
\wedge L_1\wedge\cdots\wedge L_m].
\tag{6.2}
\]
Clark calls this the general form of a clause. Suppose there are exactly \(k\)
clauses, \(k>0\), in the relation \(R\). Let
\[
\begin{array}{c}
R(x_1,\ldots,x_n)\leftarrow E_1\\
\vdots\\
R(x_1,\ldots,x_n)\leftarrow E_k
\end{array}
\]
be the \(k\) general forms of these clauses. Each of the \(E_i\) is an
existentially quantified conjunction of literals, as in (6.2). The clauses of
\(R\) can now be written in a single definition as
\[
\forall(x_1\ldots x_n)[R(x_1,\ldots,x_n)\leftrightarrow E_1\vee\cdots\vee E_k].
\tag{6.3}
\]
Taking a lead from Robinson's SUPER~[49], (6.3) is rewritten as the equation
\[
R=\lambda x_1\ldots x_n.(E_1\vee\cdots\vee E_k)
\]
which will be written in THOR-L as
\begin{alltt}
R \(\equiv\) (LAMBDA (\(x\sb{1} \cdots x\sb{n}\)) (\(\vee\) \(E\sb{1} \cdots E\sb{k}\))).
\end{alltt}
Each of the \(E_i\) general forms will be written
\begin{alltt}
(LAMBDA (\(y\sb{1} \cdots y\sb{p}\)) (\(\wedge\) (UNIFY \(z\sb{1}\) \(t\sb{1}\)) \(\cdots\) (UNIFY \(z\sb{n}\) \(t\sb{n}\))
                     \(L\sb{1} \cdots L\sb{m}\))).
\end{alltt}
The binary primitive UNIFY will be introduced in the next section; the meaning
of \(\wedge\) and \(\vee\) will be discussed in the section on search.

The list membership predicate is presented as an example of how a relation is
represented in THOR-L. The list membership predicate is defined by the two Horn
clauses
\[
\begin{array}{l}
member(x,[x|xs])\leftarrow\\
member(x,[y|ys])\leftarrow member(x,ys).
\end{array}
\]
The equivalent THOR-L definition is
\begin{alltt}
MEMBER \(\equiv\) (lambda (x list)
  (\(\vee\) (lambda (xs) (unify [x|xs] list))
     (lambda (y ys) (\(\wedge\) (unify list [y|ys])
                       (member x ys))))).
\end{alltt}

\subsection{Unification}

A strict binary operator, UNIFY, is added to THOR-L to accomplish unification.
UNIFY is a boolean valued predicate which attempts to make two
\(\lambda\)-expressions \(\alpha\)-equivalent through the binding of free
variables occurring in the expressions. The bindings made through unification
are conditional---they can be undone in the event unification fails or the
resolution search backtracks beyond the point where binding occurred. The
implementation of UNIFY is discussed in Section 6.4.

\subsection{Search}

The search strategy used by a logic programming system has a tremendous impact
on how the system is implemented, its efficiency, and its completeness. The
search strategy chosen for THOR-L is implemented by the \(\wedge\) and \(\vee\)
operators and
should fit within the constraints of the RED2 architecture. The two strategies
considered for use in THOR-L were breadth-first and depth-first.

The breadth-first strategy explores the search space in a quasi-parallel manner
that tries all possible resolvents at each step. Breadth-first systems are
complete (i.e., find all possible solutions), but are expensive in both time and
space. An explicit control mechanism is required to manage the search and it is
necessary to represent several derived clauses in memory at the same time. These
requirements do not mesh well with the existing RED2 architecture, so
breadth-first search is not used by THOR-L.

The depth-first strategy explores the search space in a serial manner,
attempting only one resolvent at each step. Backtracking is used to explore
alternative search paths. Depth-first systems are not complete, but can be
efficient in both time and space. Only one derived clause needs to be
represented at a time, which can be accomplished by storing the variable
bindings generated during unification in the redex store. Depth-first search
can be implemented on RED2 without an explicit control mechanism by letting
\(\wedge\) be the AND primitive and \(\vee\) be the OR primitive. Because
depth-first search can be
implemented on RED2 without requiring modifications, it was chosen as the THOR-L
search strategy.

The logic programming capability obtained by using the relation representation
scheme of Section 6.2.1 with AND for \(\wedge\) and OR for \(\vee\) is
equivalent to depth-first
SLD-resolution with a leftmost computation (selection) rule~[5]. Only a single
refutation is derived for a goal; to generate more than one derivation, a
different representation/compilation scheme for relations must be used. (See
Section 6.8 for an example of how to generate all refutations.)

\subsection{Interfacing Functions and Relations}

In THOR-L, relations are treated as boolean valued functions; a goal expression
may be used anywhere a boolean value is appropriate. When a goal expression
reduces to TRUE, it is a logical consequence of the relations present in the
definition context \(\delta\). The correct answer substitution generated during
the refutation of a goal expression is made available for use via the redex
store.

Functions may be used within relations in two different ways: boolean valued
functions can be used as literals within a relation, and arbitrary functional
expressions can be used as arguments to the UNIFY primitive.

\section{The Unification Algorithm}

Unification in THOR-L is different from that in conventional logic programming
languages (such as PROLOG~[16]) because THOR-L allows first-order unification of
quantified \(\lambda\)-expressions. Unification of quantified expressions is
necessary because of THOR-L's partial evaluation semantics. The primary
difference between quantified and unquantified first-order unification is that
an algorithm for the former must ensure that only free variables are bound and
that their bindings do not violate any scope restrictions. The following
examples illustrate the unification of quantified terms:

\begin{itemize}
\item The terms \((\lambda x.x)\) and \((\lambda y.y)\) unify under the identity
substitution, because the two differ only by a change of bound variable name.
\item The terms \((\lambda x.ux)\) and \((\lambda y.vy)\), where \(u\) and \(v\)
are free variables, unify under the substitution \(\{u/v\}\).
\item The terms \((\lambda x.x)\) and \((\lambda x.a)\), where \(a\) is a
constant, do not unify. Only free variables can be bound during unification.
\item The terms \((\lambda x.x)\) and \((\lambda x.u)\), where \(u\) is a free
variable, do not unify. Bound variables, such as \(x\), are local to the scope
of their quantifier and cannot be ``pulled'' out of this scoping.
\end{itemize}

Unlike higher-order unification of quantified terms, which has received much
attention~[27], the literature on first-order unification is scant, consisting
primarily of work done by Staples and Robinson [57] [55]. The Staples-Robinson
unification algorithm operates on expressions which use named identifiers to
represent variables and requires special mechanisms to avoid name clashes and to
determine if a variable is free or bound. The unification algorithm presented
here operates on expressions which use DeBruijn indices to represent variables;
this eliminates name clashes and makes the determination of free and bound
variables trivial.

The algorithm upon which the UNIFY primitive is based consists of two
procedures: Q-UNIFY and OCCURS. Unification is initiated by Q-UNIFY, which takes
as its input two \(\lambda\)-expressions to be unified and a special memory
vector, \(\sigma\), for storing bindings. If unification is successful, Q-UNIFY
returns the value true and \(\sigma\) contains the most general unifier which
makes the expressions equivalent. If unification is unsuccessful, Q-UNIFY
returns false and the contents of \(\sigma\) are undefined.

The length of \(\sigma\) is equal to the number of free variables existing at
the time of unification. Each free variable is associated with one element of
\(\sigma\), which can be accessed using the free variable's DeBruijn index. When
Q-UNIFY is called, each element of \(\sigma\) should initially be unbound,
signified by containing the value \THOR{(UBV i)} where \(i\) is the index of
the element's free variable.

\begin{alltt}
Q-UNIFY \(e\sb{0}\) \(e\sb{1}\) \(\sigma\) = U* 0 \(e\sb{0}\) \(e\sb{1}\)

where
U* \(n\) \(\lambda.e\sb{0}\)     \(\lambda.e\sb{1}\)     =  U* (\(n\)+1) \(e\sb{0}\) \(e\sb{1}\)
U* \(n\) (\(e\sb{0}\) \(e\sb{1}\))   (\(e\sb{2}\) \(e\sb{3}\))   =  (U* \(n\) \(e\sb{0}\) \(e\sb{2}\)) and (U* \(n\) \(e\sb{1}\) \(e\sb{3}\))
U* \(n\) (VAR \(i\)) (VAR \(i\)) =  true
U* \(n\) (VAR \(i\)) \(e\)       =  if \(i < n\) then false                 \(\ll\)
                              else U* \(n\) \(\sigma[i-n]\) \(e\)
U* \(n\) \(e\)       (VAR \(i\)) =  U* \(n\) (VAR \(i\)) \(e\)
U* \(n\) (UBV \(i\)) (UBV \(i\)) =  true
U* \(n\) (UBV \(i\)) \(e\)       =  if \(\sigma[i]\) = (UBV \(i\)) then
                                if (OCCURS \(i\) \(n\) \(e\) \(\sigma\)) then false
                                else \(\sigma[i]\leftarrow e\); true
                              else U* \(n\) \(\sigma[i]\) \(e\)
U* \(n\) \(e\)       (UBV \(i\)) =  U* \(n\) (UBV \(i\)) \(e\)
U* \(n\) \(e\sb{0}\)     \(e\sb{1}\)     =  false
\end{alltt}

The Q-UNIFY procedure is identical in most respects to unification procedures
for unquantified expressions; the differences being the argument \(n\) to U*
and the test performed on the line marked by \(\ll\). These two differences are
necessary to ensure Q-UNIFY only binds free variables. Free variables can easily
be detected, because their binding index will reach beyond the
\(\lambda\)-bindings encountered thus far by U*. The number of
\(\lambda\)-bindings encountered is recorded by \(n\), and the test for ensuring
only free variables get bound is performed on the line marked by \(\ll\).

The OCCURS function performs an occurs check to detect if the free variable
represented by index \(i\) occurs in the expression \(e\) and also checks to see
if binding variable \(i\) to \(e\) will cause any scope violations. If either
of these conditions is true, then OCCURS returns true; otherwise OCCURS returns
false.

\begin{alltt}
OCCURS \(i\) \(ub\) \(e\) \(\sigma\) = O* 0 \(ub\) \(e\)

where
O* \(lb\) \(ub\) \(\lambda.e\)       =  O* (\(lb\)+1) (\(ub\)+1) \(e\)
O* \(lb\) \(ub\) (\(e\sb{0}\) \(e\sb{1}\))   =  (O* \(lb\) \(ub\) \(e\sb{0}\)) or (O* \(lb\) \(ub\) \(e\sb{1}\))
O* \(lb\) \(ub\) (VAR \(j\))     =  if \(lb < j < ub\) then true
                              else if \(j \geq ub\) then O* \(lb\) \(ub\) \(\sigma[j-lb]\)
                              else false
O* \(lb\) \(ub\) (UBV \(j\))     =  if \(i = j\) then true
                              else false
\end{alltt}

The O* function uses an upper and lower bound on allowable binding indices to
detect if a scope violation occurs. The upper bound, \(ub\), is the minimum
index value for globally free variables and the lower bound, \(lb\), is the
maximum index value for variables local to \(e\). If a variable falls between
\(ub\) and \(lb\) it is locally free with respect to \(e\), but not globally free
with respect to the expressions being unified. Therefore, binding variable
\(i\) to \(e\) would cause a scope violation.

\section{Implementing Unification on RED2}

The implementation of unification is a straightforward adaptation of the
Q-UNIFY algorithm to the RED2 architecture, with one exception: the algorithm is
extended to handle lazy structures by reducing the structure components before
attempting their unification. The UNIFY primitive reduces its two arguments and
transforms the unification subgraph to fire the U* primitive, which is private
to RED2 and cannot be accessed directly from THOR. Recursion is avoided by
rewriting the U* expression according to the rules in the Q-UNIFY algorithm; for
example, the expression
\[
\THOR{(U* }n\ (e\sb{0}\ e\sb{1})\ (e\sb{2}\ e\sb{3})\THOR{)}
\]
is transformed into the expression
\[
\THOR{(AND (U* }n\ e\sb{0}\ e\sb{2}\THOR{) (U* }n\ e\sb{1}\ e\sb{3}\THOR{))}.
\]
Each rewrite costs one reduction. All of the graph rewrites can be done in situ
except for the application and structure cases. The OCCURS check is not
performed by rewrites, but is encoded directly in the RED2 control unit. In
lieu of the recursive algorithm given here, a looping algorithm is used.

The introduction of unification requires several modifications to the RED2
architecture:

\begin{itemize}
\item Indirection pointers are introduced into the redex store to bind one
variable to another.
\item A reset list is necessary to implement conditional variable bindings. The
AND, OR, and NOT primitives must be changed to handle the reset list.
\item A separate memory area, called the unification memory, is needed for
storing reduced expressions which get bound to variables.
\item References to unbound variable objects must be replaced by environment
pointers (EP instructions).
\end{itemize}

Each of these changes will now be described in detail.

\subsection{Indirection Pointers}

In the Q-UNIFY algorithm, \(\sigma\) is a contiguous vector of memory. This
characteristic allows a direct mapping from a free variable's index to its
associated memory element in \(\sigma\). Because of this direct mapping, the UBV
designator can be used to mark both unbound variables and indirection pointers
to other memory elements. Such indirection pointers are used to bind one
variable to another. Whenever \(\sigma[i]=\THOR{(UBV }j\THOR{)}\) and
\(i\ne j\), the variable \(i\) has been bound to the variable \(j\). Conversely,
\(i=j\) signals that variable \(i\) is unbound.

The RED2 equivalent of \(\sigma\) is the redex store. The currently active path
in the redex store is not generally a contiguous vector of memory, so there is
not a direct mapping from a free variable's index to its value in the redex
store. This makes it impossible to use the UBV opcode to mark both unbound
variables and indirection pointers. The UBV opcode is used only to mark unbound
variables, and a new opcode, IP, is introduced to mark indirection pointers. The
data field of an indirection pointer contains the address of another variable
in the redex store.

\subsection{Backtracking and the Reset List}

The bindings made during unification are contingent upon a successful refutation
of the goal expression. Should the current resolution search path fail to find a
refutation, some unification bindings may need to be undone by restoring the
variables' original values. To facilitate the undoing of bindings the UNIFY
primitive records, in the reset list, the address and value of each variable it
binds. The backtracking mechanism uses the reset list to undo bindings.

The reset list is actually a stack and requires a separate memory space. The
RED2 memory model can be modified so the reset list and the control stack share
adjoining memory areas and grow towards each other, like the graph memory and
the environment area do.

Backtracking in THOR-L is implemented through the combined efforts of the AND,
OR, NOT, and UNIFY primitives. Each primitive records the top of the reset list
when initially fired. If the primitive reduces to FALSE, all of the variables
bound since the initial firing (those entries in the reset list beyond the top
entry recorded by the primitive) are reset to their original values.

\subsection{Unification Memory}

Currently, RED2 has a policy of discarding subgraphs which do not become part of
the final result graph. (This policy is discussed in detail in Section 5.2.1.)
This policy is incompatible with a strict unification primitive because a
reduced expression which gets bound to a variable during unification may get
discarded before it is no longer needed. For example, in the expression
\begin{verbatim}
(LAMBDA (X Y) (IF (AND (UNIFY Y (+ X (+ 2 3))) (UNIFY X 1)) Y 0))
\end{verbatim}
the variable Y gets bound to \THOR{(+ X 5)} by UNIFY. Because the
\THOR{(+ X 5)} subgraph is not immediately part of the result graph, it is
discarded before it can be used in the true branch of the conditional.

To remedy the premature discarding of expressions, the UNIFY primitive copies
the expressions it binds into a special memory area called the unification
memory. The unification memory is managed in a stack-like fashion; upon
backtracking, it is possible to release the memory taken by expressions which
were bound to variables that have been reset.

\subsection{EP's and Delayed Variable Binding}

RED2 was initially designed to work with \(\beta\)-reduction as the only binding
mechanism. In \(\beta\)-reduction, a variable is immediately bound to a value
which cannot be changed. As a result, variables only need to dereference their
value once, during forward execution. If a variable's value is found to be an
unbound variable object, this indicates that the variable's
\(\lambda\)-abstraction is not part of a \(\beta\)-redex.

The addition of unification to RED2 changes this situation. With unification,
the value bound to a variable might not be fixed until some later time. Such a
binding is called a delayed binding. The value of a delayed binding becomes
refined as other variables get bound through unification. Because a variable
might not yet be bound, dereferencing its occurrences only during forward
execution does not always suffice. If a variable dereferences to an unbound
variable object during forward execution, the variable should also be
dereferenced during reverse execution to detect any bindings that may have
occurred between the forward and reverse execution times.

The role of environment pointers (EP instructions), which were introduced in
Section 5.2.2 to implement sharing of atomic objects, is now generalized to
optimize variable dereferencing during reverse execution. Each time a variable
instruction which is not at the head of a spine is executed in the forward
direction and dereferences to an unbound variable object, an EP instruction is
placed into the result graph with its data field containing the address of the
variable's redex store entry. Upon reverse execution, the EP instruction checks
the redex store entry to see if the variable has become bound through
unification. If it has, its value is executed; if it has not, the EP instruction
turns back into a variable instruction.

Note that this modification does not necessarily make the value of delayed
bindings available to every one of the variable's occurrences. It only makes the
value available to those occurrences which come after the binding has been made.
For example, in the expression
\begin{alltt}
(LAMBDA (X) ... X ... (UNIFY X 2) ... \underline{X} ...)
\end{alltt}
the binding of X to 2 is only seen by the underlined occurrence of X. To avoid
``missing'' a variable binding, the unifications in a relation's definition must
be carefully ordered. Such operational concerns are unfortunate because they
detract from the declarative nature of THOR-L programs.

\section{Managing Explicit Quantification}

The use of explicitly quantified variables in a language with unification
requires two problems to be solved: removal of a quantifier when its variable
has become bound, and extending a quantifier's scope when its variable becomes
part of the binding of a variable with larger scope.

The first problem is demonstrated by the expression
\begin{verbatim}
(LAMBDA (X Y) (IF (UNIFY X 1) (+ X Y) (- X Y))).
\end{verbatim}
When the X gets bound by unification, its quantifier should disappear from the
result, which should be the expression
\begin{verbatim}
(LAMBDA (Y) (+ 1 Y)).
\end{verbatim}
The second problem is more subtle. In the expression
\begin{verbatim}
(LAMBDA (X Z) ... (LAMBDA (Y) ... (UNIFY X [Y|Y]) ...) ...),
\end{verbatim}
The variable X gets bound to an expression structure containing Y, which is a
quantified variable of smaller scope. In the resulting expression, the scope of
Y should be extended to include the scope of X:
\begin{verbatim}
(LAMBDA (Y Z) ...).
\end{verbatim}

It is interesting to note that most logic programming languages do not encounter
these problems. In PROLOG, for instance, all variables are implicitly
universally quantified at the outermost level of a clause, so there are no
quantifiers to dissolve upon binding and the ordering of the quantifiers (even
if they were explicit) is not important. The use of explicit existential
quantifiers in Clark's representation of clauses makes it necessary for THOR-L
to solve these two problems.

\subsection{Removing LAMBDA's Bound Through Unification}

When a LAMBDA instruction which is not part of a \(\beta\)-redex is executed in
the forward direction, it puts an unbound variable object into the redex store
and copies itself into the result graph. If the variable associated with such a
LAMBDA instruction later becomes bound through unification, the LAMBDA
instruction must be removed from the result graph and the binding indices of
each variable in the LAMBDA's body must be decremented by one to compensate for
the lost abstraction.

To implement this behavior, the LAMBDA instruction must be modified so that when
executed in reverse it checks to see if the redex store entry associated with
the LAMBDA no longer contains an unbound variable object. To facilitate this
check, a pointer to the entry is pushed onto the control stack during forward
execution. During reverse execution the pointer is popped off the control stack
and used to access the LAMBDA's redex store entry. If the entry is no longer an
unbound variable object the LAMBDA instruction erases itself by turning into an
EX-LAMBDA instruction. The printing mechanism will ignore the EX-LAMBDA and take
care of decrementing the binding indices of the variables in the EX-LAMBDA's
body.

\subsection{Extending the Scope of Quantifiers}

The solution of the scope extension problem is quite involved, and only a brief
outline will be presented. The solution is accomplished through the cooperation
of the LAMBDA instruction, the occurs check, the print mechanism, and a new data
structure called the extension list. The extension list contains information on
every variable whose scope needs to be extended and by how much.

The first step is to discover when a variable's scope needs to be lifted. This
situation is detected during the occurs check by noting if the expression to
which a variable is being bound contains an unbound free variable of smaller
scope. If this is true, the environment entry for the unbound variable is marked
as needing to have its scope extended and the extension list records how far it
needs to be extended. If the unbound variable has already been marked for scope
extension, the extension list is updated to record the largest extension
necessary.

When a LAMBDA instruction is executed in reverse, it checks to see if its
associated redex store entry is an unbound variable which has been marked for
scope extension. If this is the case, it turns into an EX-LAMBDA instruction.

The print mechanism bears much of the responsibility for scope extension. When
an EX-LAMBDA instruction is encountered by the printer, it signifies one of two
possibilities:

\begin{enumerate}
\item The variable which was quantified by the EX-LAMBDA was bound to some
expression during unification, and this expression may contain variables whose
scope needs to be extended to this point.
\item The variable which was quantified by the EX-LAMBDA is unbound and needs to
have its scope extended.
\end{enumerate}

Only possibility (1) requires action by the print mechanism. The extension list
is searched for variables whose scope needs to be extended to this point. If any
are found, they are printed as part of a LAMBDA binding at this time. If none
are found, nothing needs to be printed. The print mechanism also handles
adjusting binding indices of variables to properly reflect the extended scopes.

\section{Adapting Unification to Incremental Reduction}

An interesting problem arises when unification is added to a system which
supports incremental reduction---how can conditional bindings be preserved
across different reduction sequences? If the expression
\begin{verbatim}
(LAMBDA (X Y) (IF (AND (UNIFY X 1) (UNIFY Y 2)) ...))
\end{verbatim}
is incrementally reduced with a quantum of 1, the \THOR{(UNIFY X 1)} reduces to
TRUE and X gets bound to 1 and no more reductions are allowed. The binding of X
to 1 is conditional upon the unification of Y with 2, so the X binding cannot
yet be allowed to extend beyond the AND expression. How can the conditional
binding be preserved until the expression is reduced again?

The BIND primitive is introduced to preserve conditional bindings across
incremental reduction sequences. When the quantum becomes exhausted during
reduction of a logical primitive (AND, OR, or NOT), all of the conditional
bindings made since the primitive was entered are recorded by inserting BIND
expressions at the point where the quantum becomes zero. For instance, the
expression from the previous paragraph reduces to
\begin{verbatim}
(LAMBDA (X Y) (IF (AND TRUE (BIND X 1) (UNIFY Y 2)) ...)).
\end{verbatim}
BIND expressions are invisible to the logical primitives; their only action is
to side-effect the redex store. They charge no reductions from the quantum and
return no result.

\section{Negation}

Logical negation in THOR-L is implemented by the NOT primitive, which uses the
\emph{negation as failure} rule~[14]. This rule states that if an attempt to prove some
goal G fails, then infer \(\neg G\). The NOT primitive reduces its single
argument \(e\); if \(e\) reduces to FALSE, \(\neg e\) is inferred to be true so
\THOR{(NOT e)} reduces to TRUE. Conversely, if \(e\) reduces to TRUE then
\(\neg e\) is inferred to be FALSE, so \THOR{(NOT e)} reduces to FALSE.

\emph{The perils of negation as failure} are well documented~[41]. The IF conditional
primitive of THOR-L eliminates the need for most uses of negation as failure and
often leads to a much clearer formulation of programs than is otherwise
possible.

The NOT primitive does not propagate any delayed bindings which may have been
made during the reduction of its argument. Therefore, double negation can be
used to test if a goal \(e\) succeeds without propagating the bindings
established during the reduction of \(e\). Smolka~[54] has suggested using
double negation to test for metalinguistic properties in functional programs.
For example, the following function returns TRUE if and only if its argument is
an unbound variable:
\begin{alltt}
VAR? \(\equiv\) (LAMBDA (X) (AND (NOT (NOT (UNIFY X 1)))
                       (NOT (NOT (UNIFY X 2)))))
\end{alltt}
Furthermore, the equality of two unbound variables can also be tested:
\begin{alltt}
EQVAR? \(\equiv\) (LAMBDA (X Y) (AND (VAR? X) (VAR? Y)
                            (NOT (AND (UNIFY X 1)
                                      (UNIFY Y 2)))))
\end{alltt}

\section{Generating All Solutions of a Goal}

It was mentioned earlier that THOR-L provides at most a single correct answer
substitution for any goal expression, even though more than one may exist. The
utility of THOR-L would be greatly enhanced if it could generate all of the
correct answer substitutions for a goal. One possible way to do this is to treat
relations as \emph{set-valued functions} which return a set or list of tuples
containing all the answers to a goal query [49] [17]. Such a solution is not
practical with the current design of RED2 because of its memory management
policies; the set of solutions can be generated, but it would be thrown away
before all the solutions could be accessed. If a separate memory area for
structures is implemented, as suggested in Section 5.2.3, the set of solutions
could safely be stored there for future use.

It is possible with the current RED2 design to generate a disjunction of tuples
containing \underline{all} of the correct answer substitutions to a goal query, albeit there
is no mechanism for accessing these tuples. This is done by using a different
clause representation and exploiting the semantics of the OR primitive.

The operational semantics given in Chapter 3 for the OR primitive specify that
non-boolean arguments are left in the disjunction as is, and reduction of the
other arguments continues. For example,
\[
\THOR{(OR 1 FALSE (+ 2 3) FALSE)}\longrightarrow\THOR{(OR 1 5)}.
\]
The OR's semantics is exploited by having each successful refutation reduce to a
tuple that contains the correct answer substitution values that are of interest.
The OR will note the value is non-boolean, will undo any conditional bindings
made during the refutation, and continue reducing the other branches in the
search.

To make this approach general, a new representation must be chosen for
relations. The new representation is based on the \emph{continuation passing} model of
PROLOG presented in~[16]. As in Clark's representation, each relation is
represented by a single function; however, the literals in each clause of the
relation are composed as nested function applications instead of being part of a
conjunction. Each function is passed its arguments and a continuation. For
example, the relation
\[
\begin{array}{l}
P(t_1,\ldots,t_n).\\
P(t_1,\ldots,t_n)\leftarrow Q(\ldots).\\
P(t_1,\ldots,t_n)\leftarrow R(\ldots),S(\ldots).
\end{array}
\]
is represented by the function
\begin{alltt}
P \(\equiv\) (LAMBDA (\(x\sb{1} \cdots x\sb{n}\) CONT)
      (OR (AND (UNIFY \(x\sb{1}\) \(t\sb{1}\)) \(\cdots\) (UNIFY \(x\sb{n}\) \(t\sb{n}\)) CONT)
          (AND (UNIFY \(x\sb{1}\) \(t\sb{1}\)) \(\cdots\) (UNIFY \(x\sb{n}\) \(t\sb{n}\)) (Q ... CONT))
          (AND (UNIFY \(x\sb{1}\) \(t\sb{1}\)) \(\cdots\) (UNIFY \(x\sb{n}\) \(t\sb{n}\)) (R ... (S ... CONT)))))
\end{alltt}
The continuation parameter CONT is used to pass an expression whose value will
be returned if the relation reduces to TRUE. This expression will be evaluated
in the environment containing the correct answer substitution.

A simple example will be given to illustrate how this scheme works. Consider the
three relations
\[
\begin{array}{l}
P(1).\\
P(2).\\
P(3).\\
Q(1).\\
Q(3).\\
P\hbox{-}intersect\hbox{-}Q(x)\leftarrow P(x),Q(x).
\end{array}
\]
The functional representation of these relations is
\begin{alltt}
P \(\equiv\) (LAMBDA (X CONT) (OR (AND (UNIFY X 1) CONT)
                         (AND (UNIFY X 2) CONT)
                         (AND (UNIFY X 3) CONT)))
Q \(\equiv\) (LAMBDA (X CONT) (OR (AND (UNIFY X 1) CONT)
                         (AND (UNIFY X 3) CONT)))

P-intersect-Q \(\equiv\) (LAMBDA (X CONT) (P X (Q X CONT)))
\end{alltt}
A disjunction of all the values for which P-intersect-Q are true can be
generated by the expression
\begin{verbatim}
(LAMBDA (Y) (P-intersect-Q Y Y)),
\end{verbatim}
which reduces to the expression \THOR{(OR 1 3)}. In this example, the goal
continuation is the expression Y. The goal continuation can be any expression
whatsoever; for example,
\begin{verbatim}
(LAMBDA (Y) (P-intersect-Q Y (+ Y Y)))
\end{verbatim}
reduces to \THOR{(OR 2 6)}.

\section{Discussion}

The experiments described here indicate that it is possible to implement some
form of logic programming capability on RED2. However, the architectural
modifications required to support unification are significant enough to indicate
that RED2 is not an ideal implementation vehicle for logic programming. In
particular, RED2's memory management policies are not compatible with delayed
binding and generating a set of all solutions of a goal. The changes
recommended in Chapter 5 to enhance sharing would also solve many of the
problems encountered in implementing THOR-L.

As a programming language, THOR-L is interesting in its own right. To my
knowledge, it is the first logic programming language with incremental reduction
semantics. THOR-L is one of the few programming languages which supports the
unification of quantified terms; QU-PROLOG~[56] and \(\lambda\)PROLOG~[40] are
two others. However, the utility of quantified unification in THOR-L is limited
by its lack of a meta-variable facility like that of QU-PROLOG which allows
unification to lift local variables out of their proper scope.

THOR-L differs from many functional languages that include unification (such as
HASL [1], FGL+LV [35], and FRESH [54]), because unification is not restricted to
pattern-directed function invocation. THOR-L allows unification to be used
anywhere a boolean value is appropriate.
